% VERSION MAP (v0.93 = advisor 9/14 items: standard class-incremental protocol, corrected utility floor, two-level selector as pre-registered, Table 1 both orders; numbers from outputs/exp3): v0.5 method-first (09-01) | v0.6 question-first + genealogy (09-03 Fable) | v0.61 narrative-first (09-03 Sol) | v0.7 = Sol spine + Fable body, PI rulings (09-03) | v0.71 = Sol red-team R2 adjudicated (09-03) | v0.8 = SOTA comparison filled (09-05) | v0.81 = PathSelect row = hier, flat as ablation, own A1/A3 rows under comparison protocol, paired 10-fold statistics; OPCM sentence corrected (09-05) | v0.82 = Sol R3 adjudicated: quoted-vs-rerun framing, no parity/SOTA wording, memory storage boundary, controlled-variants block, 2x2 novelty framing (09-06) | v0.83 = PI restructure: main table = 12 quoted + PathSelect only; ablations rebuilt as leave-out narrative; six-question and lifecycle tables removed; codes stripped (09-06) | v0.84 = forward-order results and order-dependent hierarchy reported (09-06) | v0.85 = forward-order leave-out chain; robustness synthesis (09-06) | v0.86 = QPMIL-VL reproduction check with released code, both orders (09-07); v0.87 | v0.88 = leave-out and order tables merged; float barrier before Discussion (09-07) | v0.89 = Fig.1 v2.0 integrated (09-08); v0.89a = batch-size deviation disclosed; quoting-precedent sentence made precise; benchmark named after ConSlide (09-08) | v0.89c = reverse reproduction rerun with the paper's mini-batch reproduces published values in both orders (09-08) | v0.89d = Abstract, Introduction, Contributions rewritten (Sol proposals adjudicated); attribution caveat moved to protocol (09-08) | v0.89e = Related Work rewritten (09-08) | v0.89f = Method rewritten for precision: L_sem prior defined, replay data flow, defensive parenthetical removed (09-08) | v0.89g = fixed eight-way head disclosed (09-08) | v0.89h = Sol cold-read of Method adjudicated: E, normalization, per-layer merge, softmax in KL, empty groups, hyperparameters and results moved to Sec. 4 (09-08) | v0.89i = code facts confirmed (active groups, equal-weight E, KL on softmax); seen-class check sentence (09-08) | v0.89j = Experiments restructured: labeled protocol paragraphs, question-headed ablation paragraphs (09-08) | v0.89k = Sol cold-read of Sec. 4 adjudicated: naming, scoping of attribution claims, caption shortened, Table 3 caption corrected (09-08) | v0.89l = Discussion/Limitations merged into three paragraphs with five bounded claims; Conclusion compressed (09-08) | v0.89m = official ICASSP 2027 kit files installed (byte-identical to prior); ethics statement added (09-08) | v0.89n = Sol full-manuscript consistency read adjudicated: Table 3 narrative corrected, terminology fixed (09-08) | v0.89o = figure path restored to v1_0 (09-08) | v0.89p = compression cut 1: abstract to 150 words, Related Work to two paragraphs (09-08) | v0.89q = compression cut 2: 3.3, 3.5, 4.1, 4.2 tightened, no number changes (09-08) | v0.89r = compression cut 3: Tables 3+4 merged, limitations to three points, conclusion shortened (09-08) | v0.89s = cut 4a: 5+1 layout for advisor review (09-09) | v0.89t = three-way cold read: result-independent wording fixes (STE, uniform-pooling hinge, U_old provenance, overclaims) (09-11) | v0.89u = E0 results: per-slide Delta U with CONCH logit scale, hinge activation rates, forward-order trajectory sentence, flat F_g note (09-11) | v0.89v = structural rewrite: argument-first paragraphs, info-flow Method, evidence-chain Experiments, attackable claims guarded, engineering labels removed; main table made single-column and placed as Table 1 (09-11) | v0.89v-abs = advisor's abstract adopted verbatim; guards covered in Sec. 3.3/4.2/4.3 (09-12) | v0.89v-ov = Method overview paragraph added per advisor; protocol details trimmed (09-12) | v0.89w-a = E1 zero-shot top-8 reference row and E3 selection-vs-weighting decomposition integrated (09-12) | v0.89w = E2 gated-vs-always-on control integrated; hinge described as a design choice (09-12) | v0.9 = Sol final check adjudicated (scope, attribution, terminology); advisor review base (09-12) | v0.91a purple benchmark naming (09-14) | v0.91b brown: items 3,5,7,8,9,11 (09-14) | v0.91b-c = contributions rewritten for first-impression readers (PI draft, Sol and Fable merged) (09-14) | v0.91b-c blue: hinge activation rates moved from Sec. 3.3 to 4.2 (item 4, partial) (09-14) | v0.91e green: item 10 — Related Work merged into the Introduction as a gap paragraph; contributions shortened (09-14) | v0.91e-r = authors updated, funding moved to comment, open-comment notes added (09-14) | v0.92a green: contributions enriched to ~125 words; two evidence clauses restored (09-14) | v0.92b = integrated related work (B2): two lineage paragraphs with named methods, gap→PathSelect (09-14) | v0.92c = Sol focus read: scoping (adaptation confined to selection; utility regression penalized; retention 'primarily evaluated'; 'directly targets'), contribution 1 method-heavy, CoMEL attribution generic (09-14) | v0.91d green: item 12 — results paragraphs rewritten claim→evidence→interpretation (09-14) | v0.91c blue: item 4 boundary cleanup — Sec. 3 holds definitions and rationale only; capacities, variant definitions, and E2 moved to Sec. 4 (09-14) | v0.91-g1 = benchmark naming + Sec. 4.1 definition + eight subtype classes, marked (09-14)
\documentclass{article}
\usepackage{spconf,amsmath,amssymb,graphicx,hyperref}
\usepackage{booktabs}
\usepackage{multirow}
\usepackage{placeins}
\usepackage{xcolor}
\usepackage[normalem]{ulem}
% ---- revision markup: \markuptrue = grey strike-out + blue insertions (advisor view); \markupfalse = clean ----
\newif\ifmarkup \markuptrue
\ifdefined\MARKUP \ifnum\MARKUP=1 \markuptrue\else\markupfalse\fi \fi
\newcommand{\del}[1]{\ifmarkup\textcolor{gray}{\sout{#1}}\fi}
\newcommand{\ins}[1]{\ifmarkup\textcolor{blue}{#1}\else#1\fi}
\newcommand{\rep}[2]{\del{#1}\ins{#2}}
\newcommand{\note}[1]{\textcolor{teal}{\textit{[AH] #1}}}     % PI's open comments (green)
\newcommand{\noteHF}[1]{\textcolor{red}{\textit{[HF] #1}}}    % advisor's comments (red)
\newcommand{\repc}[3]{\ifmarkup\textcolor{gray}{\sout{#2}}\textcolor{#1}{#3}\else#3\fi}
\newcommand{\repP}[2]{\repc{violet}{#1}{#2}}  % purple: v0.91a benchmark naming
\newcommand{\repB}[2]{\repc{brown}{#1}{#2}}   % brown: v0.91b items 3,5,7,8,9,11
\newcommand{\repL}[2]{\repc{blue}{#1}{#2}}    % blue: item 4
\newcommand{\repG}[2]{\repc{teal}{#1}{#2}}    % teal: items 10,12
\newcommand{\insB}[1]{\ifmarkup\textcolor{brown}{#1}\else#1\fi}
\newcommand{\delB}[1]{\ifmarkup\textcolor{gray}{\sout{#1}}\fi}
\renewcommand{\topfraction}{0.95}
\renewcommand{\bottomfraction}{0.6}
\renewcommand{\textfraction}{0.05}
\renewcommand{\floatpagefraction}{0.85}
\setcounter{topnumber}{3}
\setcounter{totalnumber}{5}
\setlength{\textfloatsep}{9pt plus 2pt minus 2pt}
\setlength{\floatsep}{8pt plus 2pt minus 2pt}
\setlength{\intextsep}{8pt plus 2pt minus 2pt}
\setlength{\abovecaptionskip}{4pt}
\setlength{\belowcaptionskip}{2pt}

\def\x{{\mathbf x}}
\def\L{{\cal L}}
\newcommand{\Fg}{F_g}
\newcommand{\Fp}{F_p}
\newcommand{\Cls}{\mathcal{C}}
\newcommand{\Psel}{\mathcal{P}}
\newcommand{\Mem}{\mathcal{M}}
\newcommand{\pending}[1]{\textbf{[#1]}}

\title{Remembering Where to Look: Continual Evidence Selection for Whole-Slide Images}

% AUTHORS — please verify (ORCID iDs are entered in the submission system, not printed):
%   Hsi-Ren Hung   ORCID 0009-0008-1857-3225  (ORCID record name: Hsi-Ren Hung)
%   Kun Lun Pan    ORCID 0009-0000-5553-944X  (ORCID record name: KUN LUN PAN; NSYSU)  alin080417@gmail.com (not public on ORCID; confirm with author)
%   Chia Feng Liao ORCID 0009-0007-5454-2277  (ORCID record name: Chia Feng Liao)  fengliao0512@gmail.com
%   Huei-Fang Yang ORCID: to be added          (corresponding author)
% FUNDING (restore at camera-ready): This work was supported in part by the National Science and Technology Council, Taiwan under Grant NSTC 112-2221-E-110-047-MY3, 114-2640-E-110-004, and 114-2634-F-006-002.
\name{Hsi-Ren Hung \qquad Kun Lun Pan \qquad Chia Feng Liao \qquad Huei-Fang Yang\sthanks{corresponding author.}}
\address{Department of Computer Science and Engineering, National Sun Yat-sen University, Taiwan}

\begin{document}
\ninept
\maketitle

% -----------------------------------------------------------------------------
\begin{abstract}
Whole-slide image (WSI) diagnosis often requires selecting a small set of diagnostically relevant patches from thousands. While continual-learning methods typically preserve predictions or representations, a budgeted system can fail differently: its diagnostic pathway may remain intact while its selector forgets where to look. \textbf{PathSelect} is a continual evidence-selection framework that learns and preserves where to look across sequential tasks. To preserve the selector, it combines replay from a bounded memory of past selections with dual-level distillation of prior selection behavior and a utility floor that penalizes loss of diagnostic usefulness without requiring the same patches to be selected. Task-local low-rank updates are folded into one shared selector that requires no task identity at inference. On the four-task TCGA benchmark under the standard class-incremental protocol, preservation raises class-incremental accuracy from 0.481 to 0.834 in the reverse task order and from 0.549 to 0.837 in the forward order. Selection-level analyses further show that retaining evidence identity is not equivalent to retaining useful evidence.
\end{abstract}

\begin{keywords}
whole-slide images, continual learning, evidence selection, computational pathology, LoRA
\end{keywords}

% -----------------------------------------------------------------------------
\section{Introduction}
\label{sec:intro}

\repG{}{A gigapixel whole-slide image (WSI) may contain tens of thousands of patches, yet practical diagnostic systems often need to identify a small subset of informative regions. When a prediction is made from a limited evidence set, missing diagnostically relevant tissue can cause failure even if the diagnostic model itself remains intact. Such systems must therefore learn where to look before the diagnosis is known, without letting the answer itself determine the evidence to be selected.}

\repG{A gigapixel whole-slide image (WSI) may contain tens of thousands of patches, yet a diagnosis often rests on a few of them. Practical systems therefore select a small evidence set and decide from it; if the selector fails to retrieve the tissue that supports the decision, the decision fails even when the classifier is sound.}{}

\repG{Learning new disease cohorts sequentially adds a specific difficulty: it changes the selector's preferences, and the evidence that old cohorts need can be crowded out. Final accuracy shows the resulting drop, but it cannot tell whether the drop originates in evidence selection or in the representation and classifier that conventional continual learning protects.}{Learning new disease cohorts sequentially adds a different failure mode: adapting to later tasks can change the selector's preferences, so that evidence needed by earlier tasks is displaced. A drop in final accuracy reveals forgetting, but not where that forgetting occurred---it may arise from the evidence-selection policy itself or from other learned components of the diagnostic pathway.}

\begin{figure*}[!t]
    \centering
    \includegraphics[width=0.8\textwidth]{figures/fig1_architecture_v1_0.pdf}
    \caption{\textbf{PathSelect overview.} Along the diagnostic path, all model components except the evidence selector (orange) are frozen. Within a task the selector allocates a budget of $B=8$ patches across tissue groups and then within groups; only the selected evidence reaches the frozen cosine head, and inference uses this solid path alone. Across tasks, a Selection Memory of past selector snapshots (scores, indices, utilities; patch features are reloaded from the training store for replay) drives replay, dual-level distillation of old scores, and a utility floor that penalizes regression while allowing the evidence to change; task-local LoRA updates are folded into one shared selector. Selection drift is evaluated separately through evidence identity (Jaccard) and utility change ($\Delta U$).}
    \label{fig:arch}
\end{figure*}

\repG{Evidence selection in WSIs has largely been studied outside continual learning. Classical multiple-instance learning (MIL) methods such as ABMIL, CLAM, and TransMIL aggregate patch-level information into a slide representation, with attention or token interactions determining which regions influence the diagnosis [ref]. More recent methods make the act of looking explicit: HistoSelect selects evidence through a tissue-aware group-to-patch hierarchy [ref], and active-perception approaches learn where to inspect next in a slide [ref]. These methods improve how evidence is acquired for a given diagnostic setting, but they are not designed to preserve that selection policy as diagnostic tasks arrive sequentially.}{\textbf{Evidence selection in WSIs has largely been studied outside continual learning.} Classical multiple-instance learning (MIL) methods such as ABMIL, CLAM, and TransMIL aggregate patch-level information into a slide representation, with attention or token interactions determining which regions influence the diagnosis \cite{ilse2018abmil,lu2021clam,shao2021transmil}. More recent methods make the act of looking explicit: HistoSelect selects evidence through a tissue-aware group-to-patch hierarchy \cite{huang2026histoselect}, and active-perception approaches learn where to inspect next in a slide \cite{chen2026adaptivepath}. These methods improve how evidence is weighted, selected, or acquired within a fixed diagnostic setting, but do not address retention of a hard evidence-selection policy as tasks arrive sequentially.}

\repG{Continual WSI learning addresses the complementary problem of retaining diagnostic competence across tasks. ConSlide rehearses stored slides [ref], MICIL studies class-incremental MIL [ref], QPMIL-VL preserves vision--language prototypes [ref], Li et al.\ distill MIL attention with a pseudo-bag memory [ref], and CoMEL and MergeSlide explore model adaptation and model merging for continual WSI learning [ref]. These approaches protect or adapt the diagnostic model, its aggregation or attention mechanism, its prototypes, or its task-specific updates, with retention primarily evaluated through downstream diagnostic performance. Even when attention is explicitly preserved, it remains part of the learned diagnostic model rather than a standalone, hard, budgeted evidence-selection policy.}{\textbf{Continual WSI learning addresses the complementary problem of retaining diagnostic competence across tasks.} Existing approaches do so through rehearsal (ConSlide \cite{huang2023conslide}), class-incremental MIL (MICIL \cite{meseguer2024micil}), vision--language prototypes (QPMIL-VL \cite{gou2025qpmil}), attention distillation with a pseudo-bag memory (Li et al.\ \cite{li2025akd}), and model adaptation or merging (CoMEL \cite{lee2025comel}, MergeSlide \cite{bui2026mergeslide}). Their continual objective protects the diagnostic model, its aggregation mechanism, its prototypes, or its task-specific updates, with retention evaluated through downstream diagnostic performance. Even when attention is preserved explicitly \cite{li2025akd}, it is not isolated as a hard, fixed-budget selector whose selected patches alone determine the diagnosis.}

\repG{This leaves a specific gap: continual WSI methods preserve what the model knows and evidence-selection methods improve where it looks, but neither line directly targets retention of a hard evidence-selection policy across sequential tasks. We address this gap with PathSelect, which restricts sequential learning to a hierarchical evidence selector while keeping the diagnostic pathway fixed. The image and text encoders of a vision--language model and the cosine diagnostic rule that connects them are kept fixed, while only a hierarchical group-to-patch selector, built on the HistoSelect substrate, is updated across tasks, confining continual adaptation to evidence selection rather than the diagnostic pathway. Across tasks the selector is preserved at two levels: distillation keeps its old scoring preferences (behavioral continuity), and a utility hinge penalizes any regression in the diagnostic value of the selected evidence while allowing the evidence itself to change (functional continuity). Changing the selected patches is therefore not treated as forgetting; losing their usefulness is.}{This leaves a specific gap: continual WSI methods preserve what the model knows and evidence-selection methods improve where it looks, but neither line directly targets retention of a hard evidence-selection policy across sequential tasks. We address this gap with PathSelect, which restricts sequential learning to a budgeted evidence selector while keeping the diagnostic pathway fixed. The vision--language encoders and the cosine diagnostic rule that connects them remain fixed; only a tissue-aware group-to-patch selector, following HistoSelect \cite{huang2026histoselect}, is adapted across tasks. Across tasks, replay maintains competence on earlier slides, dual-level distillation preserves the selector's prior scoring behavior (behavioral continuity), and a utility floor penalizes regression in the diagnostic value of the selected evidence (functional continuity). A change in the selected patches is therefore not by itself treated as forgetting; a loss of their diagnostic usefulness is.}

\repG{On the four-task TCGA continual WSI benchmark introduced by ConSlide [ref], preservation lifts unprotected sequential selection from 0.452 to 0.854 class-incremental accuracy in the reverse task order and stays within one standard deviation of the strongest published method, while the benefit of the hierarchical component is task-order dependent.}{On the four-task TCGA continual WSI benchmark introduced by ConSlide \cite{huang2023conslide}, under the standard class-incremental protocol, preservation raises class-incremental accuracy from 0.48 to 0.83 in the reverse task order and from 0.55 to 0.84 in the forward order.}

\textbf{Contributions.}
\begin{itemize}
    \item \textbf{A selector that keeps learning without the model growing.} For whole-slide diagnosis under a patch budget, PathSelect makes the evidence selector---not the diagnostic model---the component that learns each new disease cohort. Each task trains a 16.4k-parameter low-rank update that is folded into one shared selector, so model size stays fixed and inference requires no task identity.
    \item \textbf{Preserving where to look by what the evidence is worth.} PathSelect combines dual-level distillation of the selector's earlier scores with a utility floor that penalizes regression in the diagnostic value of the evidence it selects. Evidence identity and diagnostic usefulness are treated as distinct retention targets, so preservation does not require the same patches to remain selected.
    \item \textbf{A bounded memory of past selections, replayed.} A bounded Selection Memory stores selector snapshots---selected indices, scores, and utilities, but no images---and uses them for replay, distillation, and utility preservation by reloading the corresponding training features. Its capacity is fixed at 512 snapshots, so the memory does not grow with the number of tasks.
\end{itemize}

% -----------------------------------------------------------------------------
\section{Method}
\label{sec:method}

\repG{PathSelect isolates continual forgetting at the evidence selector while keeping the diagnostic pathway fixed (Fig. [ref]). A frozen pathway turns a slide into patch features and, from any set of selected patches, into a diagnosis; only the evidence selector between them is learned. Within a task, the selector learns to spend a budget of $B=8$ patches across tissue groups and then within them, supervised by the diagnosis itself, by a class-text prior, and by the measured usefulness of each candidate patch. Across tasks, the selector is preserved by replay, distillation of old scores, and a floor on the diagnostic value of its evidence, all drawn from a memory of its own past selections, and task-local LoRA updates are folded into one shared selector. At evaluation, identity and utility drift are measured separately.}{PathSelect isolates continual forgetting at the evidence selector while keeping the diagnostic pathway fixed (Fig.~\ref{fig:arch}). A frozen pathway turns a slide into patch features and, from any selected evidence set, into a diagnosis; only the evidence selector between them is adapted. Within a task, the selector learns to allocate a budget of $B=8$ patches across and within tissue groups, guided by the diagnostic loss, a class-text prior, and the counterfactual utility of each candidate patch. Across tasks, replay, score distillation, and a utility floor preserve the selector using bounded behavioral snapshots of its own past selections, while task-local low-rank updates are folded into one shared selector.}

\subsection{From a slide to a diagnosis}
\label{ssec:setup}

\repG{}{Tasks $\tau=1,\ldots,T$ arrive sequentially, each introducing new diagnostic classes; at task $\tau$, training and inference have access only to the classes $\Cls_\tau$ observed up to $\tau$, and we write $C_\tau=|\Cls_\tau|$. A slide is a set of frozen patch features $\{x_i\}_{i=1}^{N}$ from the image encoder. A fixed vocabulary of $K=8$ tissue groups (tumor, stroma, lymphocyte, necrosis, normal epithelium, vessel, adipose, background), each defined by a small ensemble of text prompts and distinct from the diagnostic classes, is encoded once by the text encoder into $\{p_k\}_{k=1}^{K}$; each patch joins the group $G_j$ with $j=\arg\max_k\cos(x_i,p_k)$, and each group is summarized by its mean $g_j$.}

\repG{Tasks $\tau=1,\ldots,T$ arrive sequentially. A slide is a set of frozen patch features $\{x_i\}_{i=1}^{N}$ from the image encoder. A fixed vocabulary of $K=8$ tissue-grouping prompts (tumor, stroma, lymphocyte, necrosis, normal epithelium, vessel, adipose, background), distinct from the diagnostic classes, is encoded once by the text encoder into $\{p_k\}_{k=1}^{K}$; each patch joins the group $G_j$ with $j=\arg\max_k\cos(x_i,p_k)$, and each group is summarized by its mean $g_j$. The group selector $\Fg$ (a two-layer MLP, like $\Fp$) scores the prototypes, $r_j=\Fg(g_j)$, and a largest-remainder allocation converts the softmax over the scores of the non-empty groups with remaining capacity into integer quotas $b_j$ with $\sum_j b_j=B=8$; the patch selector $\Fp$ scores members, $s_i=\Fp(x_i)$, and the top $b_j$ patches of each group form the evidence set $\Psel$. The slide representation is the score-weighted, $\ell_2$-normalized mean $z=\sum_{i\in\Psel}\mathrm{softmax}(s)_i\,x_i\,/\,\bigl\|\sum_{i\in\Psel}\mathrm{softmax}(s)_i\,x_i\bigr\|_2$, and the prediction is the softmax over cosine similarities between $z$ and the frozen class-text embeddings $\{t_c\}$ of the $C=8$ diagnostic subtype classes of the benchmark (Sec. [ref]), fixed from the first task. Inference is exactly this path. $B$ bounds the evidence reaching the head; all patches are still scored. A straight-through factor $\sigma(r)_j/\mathrm{sg}(\sigma(r)_j)$ on the selected members passes the diagnostic and utility-preservation gradients back to $\Fg$.}{The group selector $\Fg$ (a two-layer MLP, like $\Fp$) scores the prototypes, $r_j=\Fg(g_j)$, and a largest-remainder allocation converts the softmax over group scores into integer quotas $b_j$ with $\sum_j b_j=B=8$; the patch selector $\Fp$ scores members, $s_i=\Fp(x_i)$, and the top $b_j$ patches of each group form the evidence set $\Psel$. For diagnosis, the selected patches are score-weighted and $\ell_2$-normalized, $z=\sum_{i\in\Psel}\mathrm{softmax}(s)_i\,x_i\,/\,\bigl\|\sum_{i\in\Psel}\mathrm{softmax}(s)_i\,x_i\bigr\|_2$, and the prediction is the softmax over cosine similarities between $z$ and the frozen class-text embeddings $\{t_c\}_{c\in\Cls_\tau}$. Inference is exactly this path. $B$ bounds the evidence reaching the head; all patches are still scored. Because the quota allocation is discrete, a straight-through factor $\sigma(r)_j/\mathrm{sg}(\sigma(r)_j)$ on the selected members lets the diagnostic and utility gradients reach $\Fg$.}

\subsection{Learning to select}
\label{ssec:within}

Within a task, $\L_{\mathrm{ev}}=\L_{\mathrm{diag}}+\beta_s\L_{\mathrm{sem}}+\beta_u\L_{\mathrm{util}}$ provides three signals. The outcome signal $\L_{\mathrm{diag}}$ is the cross-entropy (CE) of the frozen head on the selected evidence set; it supervises the joint outcome but assigns credit weakly among individual candidates. The semantic signal $\L_{\mathrm{sem}}$ aligns the patch selector with a class-text prior: with $z_{ic}=\cos(x_i,t_c)/T$, the prior $p_i\propto 1-H(\mathrm{softmax}_{c\in\Cls_\tau}(z_{ic}))/\log C_\tau$ scores how discriminative each candidate appears among the classes seen so far, and the patch-score distribution over candidates is softly aligned to it with a Kullback--Leibler (KL) term. Diagnostic class text is used only as training supervision, through this prior and the frozen-head losses, and is never provided to the selector at inference. The counterfactual-utility signal $\L_{\mathrm{util}}$ asks which candidate would actually help. Let $E$ be the currently selected set; for this counterfactual evaluation only, $E$ and $E\cup\{x_i\}$ are pooled with equal weights and scored by the head on $\Cls_\tau$, so that every candidate is evaluated under the same weights. For each unselected candidate $i$ in a top-scored subset, the marginal gain
\begin{equation}
u_i
=
\L_{\mathrm{diag}}(E)
-
\L_{\mathrm{diag}}(E\cup\{x_i\})
\label{eq:ui}
\end{equation}
is positive when adding the candidate would reduce the diagnostic loss; it forms a no-gradient teacher, and $\L_{\mathrm{util}}=\mathrm{KL}(\sigma(u)\,\Vert\,\sigma(s))$, with $\sigma$ the softmax over candidates, trains the selector to rank candidates by measured usefulness. The candidate-level $u_i$ provides within-task credit assignment, whereas the set-level utility
\begin{equation}
U_\tau(\Psel)=\log C_\tau-\mathrm{CE}_\tau(\Psel,y)
\label{eq:u}
\end{equation}
is what must be preserved across tasks.

\subsection{Preserving selection across tasks}
\label{ssec:cl}

\repG{After a task is learned, PathSelect stores a bounded set of behavioral snapshots in a Selection Memory $\Mem$: an internal sample reference, the task tag, group scores $r_{\rm old}$, candidate indices with their patch scores $s_{\rm old}$ and utilities, and $U_{\rm old}$. The memory holds no image or feature tensor; replaying a snapshot reloads the slide's frozen features from the training feature store, which replay therefore requires.}{After a task is learned, PathSelect stores bounded behavioral snapshots in a Selection Memory $\Mem$: an internal sample reference, the task tag, the indices of the selected evidence, the group scores $r^{\mathrm{old}}$, and candidate indices with their patch scores $s^{\mathrm{old}}$ and candidate-level utilities. The memory record holds no image or feature tensor; replay reloads the slide's frozen features from the training feature store, which replay therefore requires. The set-level utility used by the preservation floor is not stored as a fixed target: it is recomputed at replay time under the class vocabulary available then.}

The continual objective is
\begin{align}
\L_{\mathrm{CL}}
={}&
\lambda_{\mathrm{kd}}
\left[
\mathrm{KL}\bigl(\sigma(r^{\mathrm{old}})\Vert \sigma(r^{\mathrm{new}})\bigr)
+
\mathrm{KL}\bigl(\sigma(s^{\mathrm{old}})\Vert \sigma(s^{\mathrm{new}})\bigr)
\right]
\nonumber\\
&+
\lambda_{\mathrm{eq}}
\max\bigl(0,\,U_\tau(\Psel^{\mathrm{old}})-U_\tau(\Psel^{\mathrm{new}})\bigr)
+
\lambda_r \L_{\mathrm{diag}}^{\mathrm{rep}},
\label{eq:lcl}
\end{align}
with $\sigma$ the softmax over groups and over candidates.

\repG{}{Behavioral continuity is enforced by dual-level distillation, which preserves the old soft policy---group preferences and patch-score distributions on the stored candidate indices---without requiring the hard top-$B$ set to be identical. Functional continuity is encouraged by the utility floor. The stored old evidence $\Psel^{\mathrm{old}}$ and the currently selected evidence $\Psel^{\mathrm{new}}$ are both re-evaluated with equal-weight pooling under the classes seen at replay time, so the two are compared in the same class space. The floor penalizes only a decrease in utility, so changed evidence is allowed when its diagnostic value is preserved. Replay maintains old-task competence by applying the ordinary score-weighted diagnostic loss to replayed samples under the current head.}

\repG{Behavioral continuity is enforced by $\L_{\rm KD}$: dual-level distillation preserves the old soft policy without requiring the hard top-$B$ set to be identical. Functional continuity is encouraged by $\L_{\rm eq}$: with $U=\log C-\mathrm{CE}$ the hinge equals $\max(0,\mathrm{CE}^{\rm unif}_{\rm new}-\mathrm{CE}^{\rm unif}_{\rm old})$, both terms with equal-weight pooling of the selected set, whereas the replay loss $\L_{\mathrm{diag}}^{\mathrm{rep}}$ uses the head's score-weighted pooling; the hinge penalizes utility regression on old slides and is inactive otherwise. We use this gated form to make the utility floor explicit. Task competence is maintained by replay, which applies the ordinary diagnostic loss to replayed old samples.}{}

\repG{Consolidation uses low-rank adaptation [ref], as in its continual variants [ref], only as a fixed-cost parameterization: all terms update the current task's LoRA adapters, attached to the linear layers of $\Fg$ and $\Fp$ while the base weights stay frozen; at the task boundary each residual is folded in, $W^{(\ell)}_{\tau}=W^{(\ell)}_{\tau-1}+B^{(\ell)}_\tau A^{(\ell)}_\tau$, and the adapters are reset, so model size is fixed and inference uses one shared selector.}{Consolidation uses low-rank adaptation \cite{hu2022lora}, as in its continual variants \cite{liang2024inflora,wu2025sdlora,he2025cllora,chen2026flexem}, only as a fixed-cost parameterization. The selector is randomly initialized before the first task and is never directly optimized with full-parameter gradients. For every task, including the first, a fresh rank-4 LoRA adapter (16.4k parameters) is trained while the current selector weights remain frozen during that task; at the task boundary the residual is folded into the selector, $W_{\tau}=W_{\tau-1}+B_\tau A_\tau$, and the adapters are reset, so the deployed selector size stays fixed and inference uses one shared selector.}

% -----------------------------------------------------------------------------
\section{Experiments}
\label{sec:exp}

\subsection{Protocol and measurements}
\label{ssec:protocol}

\repG{We follow the four-task TCGA continual WSI subtype-classification benchmark introduced by ConSlide [ref]---one task per cohort with two subtypes each, ESCA (ESAD, ESCC), RCC (CCRCC, PRCC; CHRCC excluded as in ConSlide), BRCA (IDC, ILC), and LUNG (LUAD, LUSC)---using the CONCH-feature release and ten patient-level folds released with QPMIL-VL [ref]. Tasks arrive in the reverse order ESCA$\rightarrow$RCC$\rightarrow$BRCA$\rightarrow$LUNG; the forward order is also reported. We instantiate the frozen encoders with CONCH [ref] (512-dimensional features). All runs use 5 epochs per task, $B=8$ (selected from a preliminary budget sweep), LoRA rank 4, $\beta_s{=}\beta_u{=}0.1$, unit $\lambda$'s, a top-256 candidate subset for the counterfactual teacher, and at most 512 memory snapshots (3.1\,MB; $\le$256 candidates each, one replayed per training step); other settings are in the released code. The comparison protocol follows the benchmark's CONCH-feature release: ten folds with the seed equal to the fold index, and average accuracy (ACC), masked ACC (task-IL), forgetting [ref], and backward transfer (BWT) [ref]. The ablation protocol fixes fold 1 and varies five seeds; comparisons are paired by seed with win counts (5/5 systematic, 4/5 directional; no $p$-values), criteria were fixed before the runs. The controlled variants are sequential fine-tuning of the selector (SeqFT), LoRA with boundary merging only, replay, replay+distillation, and full PathSelect; the flat variant used in ablations removes $\Fg$ from selection and prediction. The eight-way head is fixed across variants; class-incremental (class-IL) accuracy takes the argmax over all eight classes at every stage (restricting it to seen classes changes forgetting by at most $0.007$), task-IL restricts it to the task's own pair, and leakage is the fraction predicted into another task's classes. Selection Jaccard is the overlap between the sets selected right after learning a task and at the end; $\Delta U=U_{\rm final}-U_{\rm own}$ is the change in equal-weight evidence utility on old-task test slides.}{We follow the four-task TCGA continual WSI subtype-classification benchmark introduced by ConSlide \cite{huang2023conslide}---one task per cohort with two subtypes each, ESCA (ESAD, ESCC), RCC (CCRCC, PRCC; CHRCC excluded as in ConSlide), BRCA (IDC, ILC), and LUNG (LUAD, LUSC)---using the CONCH \cite{lu2024conch} feature release and ten patient-level folds released with QPMIL-VL \cite{gou2025qpmil}. Tasks arrive in the reverse order ESCA$\rightarrow$RCC$\rightarrow$BRCA$\rightarrow$LUNG and in the forward order LUNG$\rightarrow$BRCA$\rightarrow$RCC$\rightarrow$ESCA; class prompts are added as tasks arrive, and class-incremental accuracy is the argmax over the classes seen so far, reported after the last task. We instantiate the frozen encoders with CONCH (512-dimensional features). All runs use 5 epochs per task, $B=8$ (selected from a preliminary budget sweep), LoRA rank 4, $\beta_s{=}\beta_u{=}0.1$, unit $\lambda$'s, a top-256 candidate subset for the counterfactual teacher, and at most 512 Selection Memory snapshots, one replayed per training step; other settings are in the released code. The \emph{comparison protocol} follows the benchmark's release: ten folds with the seed equal to the fold index, and average accuracy (ACC), masked ACC (task-IL, argmax restricted to the task's own pair), forgetting \cite{chaudhry2018rwalk}, and backward transfer (BWT) \cite{lopezpaz2017gem}. The \emph{ablation protocol} fixes fold 1 and varies five seeds; comparisons are paired by seed and reported with seed-wise win counts, and the criteria were fixed before the runs. Identity and utility drift are measured separately on old-task test slides. Leakage is the fraction of old-task slides predicted into another task's classes; selection Jaccard is the overlap between the evidence selected right after learning a task and at the end; and $\Delta U$ is the change in equal-weight evidence utility between those two sets, with both scored under the final eight-class head so that the two utilities are compared in the same class space.}

\begin{table}[!t]
\centering\footnotesize\renewcommand{\arraystretch}{0.93}
\setlength{\tabcolsep}{2.0pt}
\caption{\textbf{Comparison on the four-task TCGA continual WSI benchmark} (ten-fold CV; ACC mean$_{\pm\text{sd}}$, forgetting F.\ and masked ACC Mask.\ as means). Rows above the rule are quoted from QPMIL-VL \cite{gou2025qpmil}; $^\ddagger$our rerun of its released code.}
\label{tab:sota}
\begin{tabular}{l@{\hspace{4pt}}ccc@{\hspace{7pt}}ccc}
\toprule
& \multicolumn{3}{c}{Reverse order} & \multicolumn{3}{c}{Forward order} \\
\cmidrule(lr){2-4}\cmidrule(lr){5-7}
Method & ACC$\uparrow$ & F.$\downarrow$ & Mask.$\uparrow$ & ACC$\uparrow$ & F.$\downarrow$ & Mask.$\uparrow$ \\
\midrule
Upper bound (joint) & $.908_{\pm.022}$ & --- & $.937$ & $.908_{\pm.022}$ & --- & $.937$ \\
Fine-tuning & $.234_{\pm.008}$ & $.927$ & $.803$ & $.308_{\pm.045}$ & $.841$ & $.844$ \\
EWC \cite{kirkpatrick2017ewc} & $.235_{\pm.011}$ & $.928$ & $.833$ & $.309_{\pm.051}$ & $.836$ & $.840$ \\
LwF \cite{li2017lwf} & $.236_{\pm.016}$ & $.908$ & $.900$ & $.378_{\pm.081}$ & $.731$ & $.916$ \\
A-GEM \cite{chaudhry2019agem} & $.536_{\pm.047}$ & $.527$ & $.872$ & $.436_{\pm.058}$ & $.670$ & $.879$ \\
ER-ACE \cite{caccia2022erace} & $.703_{\pm.049}$ & $.281$ & $.889$ & $.666_{\pm.049}$ & $.075$ & $.917$ \\
DER++ \cite{buzzega2020derpp} & $.684_{\pm.055}$ & $.310$ & $.910$ & $.749_{\pm.055}$ & $.219$ & $.904$ \\
ER \cite{chaudhry2019er} & $.644_{\pm.028}$ & $.387$ & $.901$ & $.790_{\pm.040}$ & $.186$ & $.894$ \\
ConSlide \cite{huang2023conslide} & $.499_{\pm.025}$ & $.058$ & $.854$ & $.659_{\pm.022}$ & $.076$ & $.861$ \\
AttriCLIP \cite{wang2023attriclip} & $.694_{\pm.058}$ & $.207$ & $.861$ & $.616_{\pm.056}$ & $.285$ & $.844$ \\
MI-Zero \cite{lu2023mizero} & $.839_{\pm.034}$ & --- & $.909$ & $.839_{\pm.034}$ & --- & $.909$ \\
QPMIL-VL \cite{gou2025qpmil} & $.859_{\pm.032}$ & $.064$ & $.925$ & $.890_{\pm.021}$ & $.027$ & $.930$ \\
\midrule
QPMIL-VL, rerun$^\ddagger$ & $.868_{\pm.046}$ & $.064$ & $.935$ & $.884_{\pm.033}$ & $.038$ & $.926$ \\
\midrule
PathSelect (ours) & $.834_{\pm.031}$ & $.101$ & $.914$ & $.837_{\pm.027}$ & $.059$ & $.916$ \\
\bottomrule
\end{tabular}
\end{table}

\subsection{Does the selector forget, and does preservation help?}
\label{ssec:ablation}

\begin{table}[!htb]
\centering\footnotesize\renewcommand{\arraystretch}{0.92}
\setlength{\tabcolsep}{3.5pt}
\caption{\textbf{Component-wise comparison} (ten-fold CV; ACC mean$_{\pm\text{sd}}$, forgetting mean). The third row omits the group level; the last row selects the eight patches with the highest class-text cosine and uses the frozen head.}
\label{tab:chain}
\begin{tabular}{lcccc}
\toprule
& \multicolumn{2}{c}{Reverse} & \multicolumn{2}{c}{Forward} \\
\cmidrule(lr){2-3}\cmidrule(lr){4-5}
Configuration & ACC$\uparrow$ & Forg.$\downarrow$ & ACC$\uparrow$ & Forg.$\downarrow$ \\
\midrule
No preservation & $.481_{\pm.057}$ & $.584$ & $.549_{\pm.099}$ & $.478$ \\
+ replay & $.819_{\pm.042}$ & $.121$ & $.831_{\pm.044}$ & $.078$ \\
+ distillation + utility floor & $.823_{\pm.041}$ & $.108$ & $.846_{\pm.047}$ & $.068$ \\
+ group-level budget (ours) & $.834_{\pm.031}$ & $.101$ & $.837_{\pm.027}$ & $.059$ \\
Class-text top-8, no learning & $.812_{\pm.024}$ & --- & $.812_{\pm.024}$ & --- \\
\bottomrule
\end{tabular}
\end{table}

\repG{Without continual preservation, the selector forgets both where to look and how useful its evidence is. Sequential fine-tuning of the flat selector keeps only $0.452$ ACC with $0.601$ forgetting in the reverse order (Table [ref]); it retains almost none of the evidence it originally chose (Jaccard $0.0023$/$0.0010$), and its old-task evidence is five to six times less useful than under full preservation ($\Delta U$ $-3.80$/$-4.71$ vs.\ $-0.74$ per slide; Table [ref]). The failure is therefore already visible at the evidence-selection level, not only as a drop in accuracy, and boundary merging alone does not address it (3/5), so consolidation by itself is not a preservation mechanism.}{\textbf{Without continual preservation, the selector forgets both where to look and how useful its evidence is.} Sequential fine-tuning of the selector keeps only $0.481$ ACC with $0.584$ forgetting in the reverse order and $0.549$/$0.478$ in the forward order (Table~\ref{tab:chain}); on fold 1 it retains almost none of the evidence it originally chose (Jaccard $0.002$) and its old-task evidence loses far more utility than under preservation ($\Delta U$ $-4.36$ vs.\ $-0.64$ or better per slide; Table~\ref{tab:ablation}). The failure is therefore already visible at the evidence-selection level, not only as a drop in accuracy.}

\repG{}{\textbf{Preservation recovers most of the loss, with replay providing the main recovery.} Replay provides the dominant recovery, raising ACC from $0.481$ to $0.819$ in the reverse order and from $0.549$ to $0.831$ in the forward order, in every fold. Distillation and the utility floor add smaller gains, to $0.823$ and $0.846$, and the group-level budget brings the final ACC to $0.834$ and $0.837$. Relative to replay alone the full model gains $0.015$ and $0.006$ ACC, improving in 6/10 and 7/10 folds; these are modest rather than fold-wise consistent improvements. In the paired fold-1 ablation the full objective exceeds replay alone by $2.5$ class-IL points in all five seeds (Table~\ref{tab:ablation}). Replay restores most old-task competence, the remaining preservation terms shape what is retained, and the group-level allocation is order-dependent.}

\begin{table}[!htb]
\centering\footnotesize\renewcommand{\arraystretch}{0.92}
\setlength{\tabcolsep}{2.2pt}
\caption{\textbf{Preservation components and what each retains} (single-level selector, fold 1, 5-seed means, reverse order). $\Delta U$ is the per-slide change in equal-weight evidence utility on old-task slides, with the initial and final evidence sets both scored under the final eight-class head.}
\label{tab:ablation}
\resizebox{\columnwidth}{!}{%
\begin{tabular}{lccc|ccc|cc}
\toprule
Arm & Rep. & Dist. & Util. & Class-IL$\uparrow$ & Task-IL$\uparrow$ & Leak.$\downarrow$ & Jacc.$\uparrow$ & $\Delta U\uparrow$ \\
\midrule
No preservation & & & & $44.5$ & $81.4$ & $48.0$ & $.002$ & $-4.36$ \\
& & \checkmark & & $57.3$ & $84.6$ & $34.8$ & $.049$ & $-2.99$ \\
& & & \checkmark & $78.9$ & $89.3$ & $13.1$ & $.070$ & $-0.88$ \\
& \checkmark & & & $77.0$ & $90.2$ & $15.0$ & $.056$ & $-0.64$ \\
& \checkmark & \checkmark & & $75.6$ & $90.6$ & $16.2$ & $.165$ & $-0.76$ \\
Full preservation & \checkmark & \checkmark & \checkmark & $\mathbf{79.5}$ & $\mathbf{90.9}$ & $\mathbf{12.7}$ & $.125$ & $\mathbf{-0.57}$ \\
\bottomrule
\end{tabular}}
\end{table}

\repG{Preservation recovers most of the loss, with replay providing the main recovery. In the leave-out chain, replay raises ACC from $0.452$ to $0.821$ in the reverse order and from $0.529$ to $0.841$ in the forward order, in every fold. Adding distillation and utility preservation further reduces forgetting in both orders ($2.4$ points, 8/10 folds, and $0.9$ points, 8/10), although the corresponding accuracy gains are modest ($0.8$ points, 7/10, and none). The hierarchical selector provides an additional benefit in the reverse order, reaching $0.854$ ACC with $0.060$ forgetting, but does not improve the forward order; this separates the more consistent contribution of preservation from the order-dependent contribution of the hierarchy (Sec. [ref]).}{}

\repG{The preservation signals address different failure modes rather than acting as interchangeable regularizers. Distillation alone stabilizes how the selector scores but not which task its evidence belongs to (leakage $0.323$; Table [ref]); replay restores old-task competence (the best single-term task-IL, $0.907$); and utility preservation gives the best single-term class-IL ($0.809$); under the hierarchy, removing the group level of distillation costs $3.95$ task-IL points (5/5), so both levels matter. Among the evaluated configurations no single loss or pair matches the full set.}{\textbf{The preservation signals address different failure modes rather than acting as interchangeable regularizers.} On fold 1, distillation alone stabilizes how the selector scores but not which task its evidence belongs to (leakage $34.8$; Table~\ref{tab:ablation}), whereas utility preservation and replay each recover most of the lost competence ($78.9$ and $77.0$ class-IL) and the full objective reaches the highest class-IL at $79.5$. \textbf{Keeping the same patches is not the same as keeping useful evidence.} The configuration with replay and distillation retains the most of its original evidence (Jaccard $0.165$), yet its accuracy stays below that of the full objective, whose Jaccard is lower ($0.125$). Evidence identity and diagnostic utility therefore capture different aspects of selector forgetting and must be read separately.}

\subsection{Where does PathSelect stand among published methods?}
\label{ssec:sota}

\repG{}{\textbf{Among learned continual methods PathSelect is second only to QPMIL-VL, and it is level with the strongest zero-shot reference.} PathSelect trails QPMIL-VL, the strongest continual-learning method in Table~\ref{tab:sota}, by $0.025$ ACC in the reverse order and $0.053$ in the forward order, and outperforms the remaining continual baselines in both orders; rerunning QPMIL-VL with its released code reproduces its published results to within $0.01$ in both orders. Among zero-shot references, MI-Zero reaches $0.839$ ACC, essentially level with PathSelect at $0.834$ and $0.837$ given the fold-to-fold spread, and a simpler class-text top-8 selector reaches $0.812$ without learning (Table~\ref{tab:chain}). Class-text similarity is therefore already a strong evidence heuristic and PathSelect does not currently surpass the best zero-shot reference. The distinction is not higher accuracy on this benchmark but learnability: PathSelect updates its evidence policy from task feedback and preserves that policy across sequential tasks, whereas the zero-shot references remain fixed.}

\subsection{Role of the group level and warm-start training}
\label{ssec:order}

\repG{}{Removing the group level, so that the selector scores all patches and selects the global top eight, changes ACC by $-0.011$ in the reverse order and $+0.008$ in the forward order while slightly increasing forgetting in both (Table~\ref{tab:chain}). We therefore retain the group level as a tissue-aware mechanism for allocating the evidence budget, but do not interpret it as a consistent source of accuracy gain. Full-parameter training of the first task likewise showed no systematic accuracy benefit in the paired fold-1 experiment, so all tasks use the same adapter-and-fold-in update.}

\FloatBarrier
% -----------------------------------------------------------------------------
\section{Discussion and Limitations}
\label{sec:discussion}

\repG{}{Replay provides most of the recovery, while distillation and the utility floor influence what is retained---selection behavior or evidence usefulness. On this benchmark, zero-shot class-text selection is already strong, and PathSelect does not outperform it in final accuracy. The value of a learned selector is therefore not established by higher accuracy here, but by its ability to adapt its evidence policy from task feedback and to preserve that policy across tasks. Whether this advantage becomes more important when class-text similarity is a weaker cue remains an open question.}

\repG{Keeping the same patches is not the same as keeping useful evidence. Distillation without the hinge retains slightly more of the original evidence (Jaccard $0.1752$ vs.\ $0.1613$) yet is less accurate and retains less utility, so preserving where the selector looks is not sufficient to preserve what its evidence contributes to diagnosis. The dominant source of the measured drift is evidence replacement rather than reweighting: on fold 1, re-scoring the old evidence set with the final selector changes the equal-weight loss by only $-0.003{\pm}0.012$ per slide, whereas replacing it with the final evidence set changes it by $+0.74{\pm}0.31$.}{}

\repG{The gating mechanism itself provides no measured advantage. In the reverse-order flat setting the hinge is active in only 1.7--7.4\% of replay steps, and replacing it with an always-on equal-weight cross-entropy on the same replayed slides changes no metric systematically across five seeds (class-IL $+0.3$, task-IL $-1.5$ points). Gating is therefore a design choice that makes the utility floor explicit, not the source of the preservation gain.}{}

\repG{PathSelect is competitive with the strongest published method in the reverse order, but this comparison does not hold in the forward order. In the reverse order it reaches $0.854{\pm}.044$ ACC and $0.060{\pm}.041$ forgetting (Table [ref]), within one standard deviation of the strongest published result ($0.859$). Rerunning that method with its released code reproduces its published results in both orders to within $0.01$ (reverse $0.868{\pm}.046$, forward $0.884{\pm}.033$); PathSelect is $0.005$ and $0.014$ below the published and reproduced reverse values with comparable forgetting ($0.060$ vs.\ $0.064$), whereas in the forward order it reaches $0.811$ (hierarchical) and $0.841$ (flat) against a reproduced $0.884$. The result therefore supports a competitive reverse-order comparison, not an overall claim of superiority. Learned selection adds value beyond a strong class-text retrieval reference, but the gain is also order-dependent. A zero-shot reference that takes the eight patches with the highest class-text cosine reaches $0.812$ without learning (Table [ref]); PathSelect is $0.042$ above it in the reverse order (9/10 folds) and level with it in the forward order, and, unlike that reference, it selects evidence without class prompts. It also adapts only 16.4k selector parameters per task (vs.\ $0.365$M reported; context only).}{}

\repG{The hierarchical substrate is not uniformly beneficial across task orders. It adds $2.5$ ACC points in the reverse order but costs $3.0$ in the forward order, where it falls behind the flat selector from the first added task (BRCA, $-0.037$ class-IL, 4/10 folds) and does not recover; we therefore treat it as an order-dependent substrate choice, not the source of the robust gain. Update consolidation alone is also insufficient to reproduce the full preservation benefit. Composing independently trained per-task residuals post hoc beats the bare sequential selector in every seed (cf.\ task arithmetic [ref]) yet trails the full method by $13.1$ points (5/5), and MergeSlide's orthogonal-projection merging [ref] does not close that gap ($0.709{\pm}.035$); the benefit comes from consolidation together with preservation, not from merging alone.}{}

\repG{The claims are bounded in three ways. First, four organ-separated, highly separable TCGA tasks, one frozen foundation model, one budget, and known task boundaries limit generality; fold-to-fold spread (sd $0.03$--$0.05$) exceeds the gaps between the top methods. Second, the hierarchy's advantage and part of the preservation gains are order-dependent. Third, the head assumes the full class vocabulary as text from the start, and replay requires the training feature store. PathSelect is a one-shot acquisition policy; sequential, stateful acquisition is left to future work.}{The claims are bounded in three ways. First, the benchmark contains four organ-separated, highly separable TCGA tasks, uses one frozen foundation model and one evidence budget, and assumes known task boundaries during training; fold-to-fold variation (sd $0.03$--$0.05$) is larger than the gaps among the top methods. Second, the additional accuracy gain beyond replay is small in the ten-fold means, with the clearest paired improvement observed only in the fold-1 five-seed ablation. Third, replay requires access to the training feature store, and no privacy guarantee is claimed. PathSelect is also a one-shot evidence-acquisition policy; sequential, stateful acquisition is left to future work.}

% -----------------------------------------------------------------------------
\section{Conclusion}

\repG{Forgetting in budgeted WSI diagnosis can occur in the evidence-selection step alone. PathSelect isolates that step behind a frozen diagnostic pathway, preserves one shared selector by keeping its old scoring preferences and penalizing utility loss; on the TCGA continual WSI benchmark it recovers most of the unprotected loss and matches the strongest published method within one standard deviation (reverse order).}{Forgetting in budgeted WSI diagnosis can occur at the evidence-selection step even when the diagnostic pathway is fixed. PathSelect isolates this failure mode and preserves one shared selector across tasks through replay, score distillation, and utility preservation, recovering most of the loss caused by unprotected sequential learning in both task orders. These results show that continual learning for budgeted perception concerns not only retaining what a model predicts, but also preserving the policy by which it acquires the evidence used to predict.}

\bibliographystyle{IEEEbib}
\bibliography{refs}

\section*{Compliance with Ethical Standards}
This study uses publicly available, de-identified whole-slide images from The Cancer Genome Atlas (TCGA), released under its data-use policies; no new human-subject data were collected and no additional ethical approval was required.

\end{document}
